\section{Related Work}\label{related-work}

Methods for fitting multiple geometric structures differ primarily in how they couple model estimation with point grouping. Consensus methods search for hypotheses supported by many observations, preference-based methods group points according to their agreement with sampled models, and energy-based methods optimize an explicit objective over assignments and model parameters. These alternatives provide the context for comparing sequential RANSAC with PEARL.

RANSAC \cite{fischler1981ransac} established the hypothesis-and-verification approach to robust estimation, in which a model generated from a minimal sample is assessed by its consensus set. Subsequent work has examined how refinement can improve the reliability of sampled estimates. Chum, Matas, and Kittler \cite{chum2003locally} observed that even an outlier-free sample need not produce a model consistent with all true inliers under measurement noise. Their locally optimized RANSAC (LO-RANSAC) refines promising hypotheses to improve support. This demonstrates the importance of distinguishing a sampled hypothesis from its refined estimate. However, improving one model does not itself resolve the competition among several structures. Sequential extraction still commits points to earlier models before later surfaces are considered. The least-squares refinement used in our baseline follows this general motivation but is not an implementation of the complete LO-RANSAC algorithm.

Preference-based clustering addresses multiple structures by representing how each observation agrees with a collection of hypotheses. Toldo and Fusiello \cite{toldo2008jlinkage} introduced J-linkage, which describes a point by the set of sampled models consistent with it and groups points using agglomerative clustering in this preference space. Magri and Fusiello \cite{magri2014tlinkage} extended this idea through T-linkage, replacing binary preferences with a continuous representation. Their distinction is relevant near uncertain inlier boundaries, where a hard threshold discards differences in residual magnitude. These methods group observations through shared model preferences rather than repeatedly removing the strongest consensus set. Their central representation is nevertheless different from the explicit spatially regularized assignment energy studied here.

Energy-based fitting makes the relationship between data association and model estimation explicit. Isack and Boykov \cite{isack2012energy} formulate multi-model fitting through geometric error and regularization, and introduce PEARL to alternate assignment optimization with model re-estimation. For a fixed candidate set, the assignment stage uses the \(\alpha\)-expansion framework of Boykov, Veksler, and Zabih \cite{boykov2001graphcuts}, allowing many observations to change labels within a move. The continuous fitting stage then updates the geometric parameters using the current groups. This combination is important because neither a fixed proposal pool nor a single pass of point assignment fully addresses their mutual dependence. PEARL provides an iterative strategy for reducing the same objective across both updates, while retaining dependence on initialization and local solutions.

The choice of regularization determines which properties an energy favors. Spatial smoothness penalizes neighboring observations that receive different labels, whereas a label cost penalizes the use of models themselves. Delong et al. \cite{delong2012labelcosts} developed an extension of \(\alpha\)-expansion that accommodates such label costs, providing a mechanism for discouraging unnecessary model complexity. These terms should not be conflated: coherent point groups do not imply that the number of models is explicitly penalized. Our implementation adopts the basic geometric-error and neighborhood formulation of PEARL, without an additional model-count term. Consequently, its plane-count cap is a computational setting rather than an optimized complexity penalty.

The computational cost of discrete assignment optimization has also motivated continuous alternatives. Amayo et al. \cite{amayo2018coral} proposed CORAL, which relaxes discrete point memberships into soft assignments and optimizes a global fitting energy using a formulation suitable for parallel computation. Their evaluation includes plane extraction from RGB-D data and homography estimation. This work shows that the energy formulation and the numerical solver are distinct design choices: similar geometric objectives can be approached through different assignment representations. It also motivates examining runtime alongside segmentation quality. The present study retains graph-cut-based expansion, while CORAL illustrates an alternative approach to assignment optimization.

Plane extraction is also shaped by the sensor model and the intended application. Kaess \cite{kaess2015planes} demonstrates the use of plane landmarks in simultaneous localization and mapping, while Schaefer et al. \cite{schaefer2019planes} estimate finite planes through a probabilistic treatment of laser measurements that incorporates ray information. The latter work also introduces SynPEB for quantitative evaluation. Our experiments use the supplied point coordinates and reference plane memberships, without adopting that measurement-likelihood formulation. Within this scope, we compare greedy extraction, sampled-proposal PEARL, and RANSAC-initialized PEARL to examine the effects of joint reassignment, initialization, and spatial regularization. J-linkage, T-linkage, LO-RANSAC, label-cost extensions, and CORAL provide methodological context rather than additional experimental baselines.
